% Standalone TikZ figure: SVM maximum-margin hyperplane.
% \input from the chapter's raw-LaTeX block; not rendered by Obsidian.
\begin{figure}[H]
\centering
\begin{tikzpicture}[
    scale=0.95,
    classA/.style={circle, draw=mondrianblue, fill=mondrianblue!70, minimum size=7pt, inner sep=0pt},
    classB/.style={rectangle, draw=mondrianred, fill=mondrianred!70, minimum size=7pt, inner sep=0pt},
    sv/.style={draw=black, thick, minimum size=12pt},
    >=stealth
]
    % Decision boundary y = 3.5 - x/7, and margins offset by +-1 in y
    \draw[very thick] (0,3.5) -- (7,2.5) node[pos=1.0, right] {\small $w^\top x + b = 0$};
    \draw[dashed, thick, mondrianblue] (0,4.5) -- (7,3.5);
    \draw[dashed, thick, mondrianred]  (0,2.5) -- (7,1.5);

    % Margin width indicator (endpoints sit exactly on the two margin lines at x=3.5);
    % the label sits above the decision boundary (which crosses this arrow at y=3.0),
    % not on top of it.
    \draw[<->] (3.5,4.0) -- (3.5,2.0) node[pos=0.25, right, xshift=2pt] {\small margin};

    % Class +1 points (blue circles), above the upper margin
    \node[classA] at (0.5,4.9) {};
    \node[classA] at (2.5,4.7) {};
    \node[classA] at (4.0,4.2) {};
    \node[classA] at (6.0,4.0) {};

    % Class -1 points (red squares), below the lower margin
    \node[classB] at (0.7,1.8) {};
    \node[classB] at (2.0,1.6) {};
    \node[classB] at (4.5,1.3) {};
    \node[classB] at (6.3,1.0) {};

    % Support vectors: sit exactly on a margin line, outlined in black
    \node[classA, sv] at (1.5,4.286) {};
    \node[classA, sv] at (5.0,3.786) {};
    \node[classB, sv] at (3.0,2.071) {};

    % Legend
    \node[classA, label=right:{\small Class $+1$}] at (8.6,4.6) {};
    \node[classB, label=right:{\small Class $-1$}] at (8.6,3.9) {};
    \node[classA, sv, label=right:{\small support vector}] at (8.6,3.2) {};
\end{tikzpicture}
\caption{The maximum-margin hyperplane of a Support Vector Machine. The decision boundary $w^\top x + b = 0$ is placed to maximize the margin between the two classes; the support vectors (outlined) are the points that lie exactly on the margin boundary and alone determine the solution.}
\label{fig:svm}
\end{figure}
